\begin{abstract}
A fiscalização convencional de licitações analisa certames e empenhos isoladamente, falhando em detectar vínculos societários, familiares e cadastrais ocultos entre concorrentes formais. Contudo, bases municipais brasileiras omitem identificadores nominais de sócios, mascaram CPFs e trazem chaves estrangeiras ausentes entre licitações e empenhos, além de exigirem proteção rigorosa de dados pessoais na comunicação científica. Este trabalho propõe um \textit{pipeline} reprodutível em Python para triagem societária preventiva e materialidade contratual, integrando oito tabelas do SICOM/TCE-MG (Divinópolis/MG, 2025) ao QSA da Receita Federal via BrasilAPI (100\% dos 673 CNPJs válidos). O método modela um multigrafo heterogêneo multicamada (\textit{Órgão}, \textit{Empresa} e \textit{Sócio Real}), pontua pares co-disputantes por evidências objetivas e onomásticas, executa redução determinística em listas de adjacência (indução, $1$-core via BFS, Union-Find e triângulos $K_3$), pondera o risco pelos R\$\,140,29 milhões empenhados líquidos do exercício e aplica pseudonimização conforme a LGPD. Como resultados, 1.788 vértices e 5.914 arestas da macro-rede ($Q=0{,}7673$) tiveram os vértices comprimidos em 98,5\% (para 26 vértices e 23 arestas em 6 componentes), isolando 3 triângulos $K_3$ e 3 núcleos críticos contratados (incluindo indício de \textit{cover bidding} em 59 itens), além de rebaixar um falso positivo onomástico de R\$\,1,36 milhão para revisão manual, executando ponta a ponta em 10,607\,s em memória.
\end{abstract}

\begin{IEEEkeywords}
Teoria dos Grafos, Grafos Multicamada, Contratações Públicas, Triângulos $K_3$, QSA, SICOM, LGPD.
\end{IEEEkeywords}
